\documentclass[12pt,a4paper]{article}

\usepackage{fontspec}
\usepackage{polyglossia}
\usepackage{amsmath}
\usepackage{amssymb}
\usepackage{geometry}
\usepackage{hyperref}

\setmainlanguage{russian}
\setotherlanguage{english}
\setmainfont{Arial}
\setsansfont{Arial}
\setmonofont{Menlo}

\geometry{margin=2.2cm}

\title{Изменение математики L2: от разомкнутой кинематики к обратной связи}
\author{RaspTank}
\date{}

\newcommand{\vdes}{v_{\mathrm{des}}}
\newcommand{\wdes}{\omega_{\mathrm{des}}}
\newcommand{\wgyro}{\omega_{\mathrm{gyro}}}
\newcommand{\thetahat}{\hat{\theta}}
\newcommand{\thetaref}{\theta_{\mathrm{ref}}}
\newcommand{\sat}{\operatorname{sat}}

\begin{document}

\maketitle

\section{Идея уровней}

В системе используются три уровня:

\[
L3 \longrightarrow L2 \longrightarrow L1.
\]

Их смысл:

\begin{itemize}
  \item \(L3\) решает, куда и как нужно двигаться.
  \item \(L2\) переводит желаемое движение корпуса в команды левой и правой
  гусеницы.
  \item \(L1\) отправляет команды на моторы и читает датчики.
\end{itemize}

В исходной постановке \(L2\) получает две величины:

\[
(\vdes, \wdes),
\]

где:

\begin{itemize}
  \item \(\vdes\) --- желаемая линейная скорость корпуса;
  \item \(\wdes\) --- желаемая угловая скорость корпуса.
\end{itemize}

На выходе \(L2\) должен получить команды:

\[
(U_L, U_R),
\]

где:

\begin{itemize}
  \item \(U_L\) --- команда левой гусеницы в процентах;
  \item \(U_R\) --- команда правой гусеницы в процентах.
\end{itemize}

\section{Что было до изменения}

\subsection{Базовая кинематика}

До изменения \(L2\) работал как разомкнутый математический преобразователь.
Он не проверял, что машинка реально едет так, как было рассчитано. Схема была
такой:

\[
(\vdes,\wdes)
\longrightarrow
(U_L^0,U_R^0)
\longrightarrow
\text{моторы}.
\]

Индекс \(0\) означает ``базовая команда без обратной связи''.

Для дифференциальной машинки используются стандартные формулы:

\[
v_R^* = \vdes + \wdes \frac{l}{2},
\qquad
v_L^* = \vdes - \wdes \frac{l}{2}.
\]

Здесь:

\begin{itemize}
  \item \(v_R^*\) --- требуемая скорость правой гусеницы;
  \item \(v_L^*\) --- требуемая скорость левой гусеницы;
  \item \(l\) --- расстояние между центрами левой и правой гусеницы;
  \item \(\wdes\) в этих формулах используется в радианах в секунду.
\end{itemize}

После этого скорости гусениц переводятся в проценты:

\[
U_R^0 = 100 \frac{v_R^*}{V_{R,\max}},
\qquad
U_L^0 = 100 \frac{v_L^*}{V_{L,\max}}.
\]

Здесь \(V_{R,\max}\) и \(V_{L,\max}\) --- скорости правой и левой гусеницы
при команде \(100\%\).

\subsection{Пример: движение прямо}

Для движения прямо:

\[
\wdes = 0.
\]

Тогда:

\[
v_R^* = \vdes,
\qquad
v_L^* = \vdes.
\]

То есть математически обе гусеницы должны двигаться с одинаковой скоростью.
Если калибровка идеальная и гусеницы ведут себя одинаково, машинка едет
прямо.

\subsection{Пример: поворот на месте}

Для поворота на месте:

\[
\vdes = 0,
\qquad
\wdes \ne 0.
\]

Тогда:

\[
v_R^* = \wdes \frac{l}{2},
\qquad
v_L^* = -\wdes \frac{l}{2}.
\]

То есть одна гусеница движется вперед, а другая назад.

\section{Почему этого недостаточно}

Реальная машинка не является идеальной математической моделью. Даже если
формулы правильные, фактическое движение может отличаться.

Основные причины:

\begin{itemize}
  \item правая и левая гусеницы имеют разную силу;
  \item гусеницы проскальзывают;
  \item трение зависит от пола;
  \item батарея проседает;
  \item при разных процентах моторов поведение меняется не строго линейно.
\end{itemize}

Поэтому одной формулы

\[
(\vdes,\wdes) \longrightarrow (U_L^0,U_R^0)
\]

недостаточно. Нужно смотреть, что происходит с машинкой фактически, и
подправлять команды.

Для этого используется гироскоп. Он дает фактическую угловую скорость:

\[
\wgyro.
\]

Также из гироскопа можно получить оценку текущего угла корпуса:

\[
\thetahat.
\]

\section{Что добавляется}

Базовая кинематика не удаляется. Она остается первым шагом:

\[
(\vdes,\wdes)
\longrightarrow
(U_L^0,U_R^0).
\]

После этого добавляется поправка:

\[
\Delta U.
\]

Итоговые команды становятся:

\[
U_L = U_L^0 - \Delta U,
\qquad
U_R = U_R^0 + \Delta U.
\]

Это очень важная часть:

\begin{itemize}
  \item если \(\Delta U > 0\), левая гусеница получает меньше, правая больше;
  \item если \(\Delta U < 0\), левая гусеница получает больше, правая меньше.
\end{itemize}

Так \(L2\) может компенсировать увод без изменения базовой кинематики.

\section{Движение вперед и назад}

Для движения вперед или назад мы хотим, чтобы машинка не поворачивала.
Значит целевая угловая скорость равна нулю:

\[
\wdes = 0.
\]

Но этого мало. Если машинка уже немного повернулась, простое удержание
нулевой угловой скорости только остановит дальнейший поворот. Оно не вернет
корпус к исходному углу.

Поэтому для движения вперед и назад используются две ошибки:

\begin{itemize}
  \item ошибка угловой скорости;
  \item ошибка угла.
\end{itemize}

\subsection{Ошибка угловой скорости}

Ошибка угловой скорости:

\[
e_\omega = \wdes - \wgyro.
\]

Для движения вперед и назад:

\[
e_\omega = -\wgyro.
\]

Если машинка поворачивает вправо или влево, гироскоп это видит, и
\(e_\omega\) становится ненулевой.

\subsection{Ошибка угла}

В начале движения запоминается начальный угол:

\[
\thetaref = \thetahat(t_0).
\]

В процессе движения есть текущий угол:

\[
\thetahat(t).
\]

Ошибка угла:

\[
e_\theta =
\operatorname{wrap}_{[-180^\circ,180^\circ]}
\left(\thetaref - \thetahat(t)\right).
\]

Функция \(\operatorname{wrap}\) нужна, чтобы угол всегда был в диапазоне
от \(-180^\circ\) до \(180^\circ\):

\[
\operatorname{wrap}_{[-180^\circ,180^\circ]}(\alpha)
=
\left((\alpha + 180^\circ) \bmod 360^\circ\right) - 180^\circ.
\]

Смысл простой: если корпус ушел от начального направления, \(e_\theta\)
показывает, на сколько градусов его нужно вернуть.

\subsection{Постоянная поправка}

Если одна гусеница стабильно слабее другой, нужна постоянная поправка:

\[
U_{\mathrm{trim}}.
\]

Она не зависит от мгновенного значения гироскопа. Это базовая компенсация
асимметрии моторов.

В общем виде она может зависеть от мощности:

\[
U_{\mathrm{trim}} = f(|U^0|).
\]

Это нужно потому, что асимметрия на малой и большой скорости может быть
разной.

\subsection{Интегральная поправка}

Если маленькая ошибка сохраняется долго, вводится накопленная ошибка:

\[
I_k =
\sat_{[-I_{\max},I_{\max}]}
\left(
I_{k-1} + \left(e_{\omega,k} + K_\theta e_{\theta,k}\right)\Delta t
\right).
\]

Она помогает убрать постоянный остаточный увод.

\subsection{Итоговая формула для вперед и назад}

Для движения вперед и назад поправка имеет вид:

\[
\Delta U =
\sat_{[-U_{\max}^{corr},U_{\max}^{corr}]}
\left(
U_{\mathrm{trim}}
+ K_\omega e_\omega
+ K_\theta e_\theta
+ K_I I
\right).
\]

Здесь:

\begin{itemize}
  \item \(U_{\mathrm{trim}}\) --- постоянная компенсация разницы гусениц;
  \item \(K_\omega e_\omega\) --- реакция на текущую угловую скорость;
  \item \(K_\theta e_\theta\) --- возврат к начальному курсу;
  \item \(K_I I\) --- компенсация долго сохраняющейся ошибки;
  \item \(\sat\) ограничивает поправку, чтобы команды не стали слишком
  большими.
\end{itemize}

\section{Поворот на месте}

Для поворота на месте задача другая. Нам не нужно удерживать начальный
курс, потому что корпус должен поворачиваться.

Условия режима:

\[
\vdes \approx 0,
\qquad
|\wdes| > 0.
\]

В этом режиме ошибка угла к начальному направлению не используется:

\[
e_\theta = 0.
\]

Используется только ошибка угловой скорости:

\[
e_\omega = \wdes - \wgyro.
\]

Поправка для поворота:

\[
\Delta U_{\mathrm{turn}} =
\sat_{[-U_{\max}^{turn},U_{\max}^{turn}]}
\left(
K_\omega^{turn} e_\omega
\right).
\]

Итоговые команды:

\[
U_L = U_L^0 - \Delta U_{\mathrm{turn}},
\qquad
U_R = U_R^0 + \Delta U_{\mathrm{turn}}.
\]

Постоянная поправка \(U_{\mathrm{trim}}\), которая нужна для движения
вперед или назад, не применяется автоматически к повороту. Иначе поворот
может стать несимметричным.

\section{Итог}

До изменения \(L2\) делал только это:

\[
(\vdes,\wdes)
\longrightarrow
(U_L^0,U_R^0).
\]

После изменения \(L2\) делает два шага:

\[
(\vdes,\wdes)
\longrightarrow
(U_L^0,U_R^0)
\]

\[
(U_L^0,U_R^0,\wgyro,\thetahat,\thetaref)
\longrightarrow
(U_L,U_R).
\]

То есть исходная математика не заменяется. Добавляется второй слой:
обратная связь по гироскопу.

Для движения вперед и назад обратная связь удерживает начальный курс:

\[
\Delta U =
U_{\mathrm{trim}}
+ K_\omega e_\omega
+ K_\theta e_\theta
+ K_I I.
\]

Для поворота на месте обратная связь удерживает заданную угловую скорость:

\[
\Delta U_{\mathrm{turn}} =
K_\omega^{turn} e_\omega.
\]

\end{document}
